\subsection{Анализ предметной области}

Городская среда представляет собой сложную систему, в которой жители ежедневно взаимодействуют с объектами инфраструктуры и муниципальными службами. Город Киров~--- административный центр Кировской области с населением около 530~тысяч человек~--- включает четыре городских района (Ленинский, Октябрьский, Первомайский, Нововятский), каждый из которых обладает собственной инфраструктурой жилищно-коммунального хозяйства, объектами благоустройства и сетью ответственных городских служб.

Городская инфраструктура Кирова охватывает широкий спектр объектов. К основным категориям относятся:
\begin{itemize}
    \item дорожное покрытие и тротуары;
    \item уличное освещение;
    \item контейнерные площадки и объекты санитарной очистки;
    \item дворовые территории и объекты благоустройства.
\end{itemize}

Каждая категория объектов закреплена за определённой организацией или ведомством:
Каждая категория объектов закреплена за определённой организацией или ведомством. Обслуживание инфраструктуры распределено между муниципальными учреждениями, унитарными предприятиями, региональными операторами и управляющими компаниями.

Ключевой проблемой взаимодействия жителя с городской средой является отсутствие единого удобного канала для сообщения о дефектах инфраструктуры и контроля хода их устранения. Обслуживание городской инфраструктуры распределено между множеством организаций: за дорожное покрытие, уличное освещение, контейнерные площадки и иные объекты отвечают разные службы, каждая из которых имеет собственные каналы приёма обращений. Для жителя это означает необходимость самостоятельно определить, какая именно служба отвечает за устранение конкретной проблемы, найти её контактные данные и составить обращение по телефону или электронной почте. Контактная информация муниципальных служб рассредоточена по разделам сайта администрации города, ведомственным порталам и справочным ресурсам, актуальность найденных данных не гарантирована.

Не менее существенные проблемы испытывают муниципальные службы: диспетчеры не располагают единым инструментом для регистрации и ведения обращений, отсутствует визуализация проблем на карте города, определение ответственной службы не автоматизировано, журнал действий оператора не ведётся. Ситуацию усугубляет отсутствие механизма обнаружения дубликатов: жители не видят обращения друг друга, поэтому одна и та же проблема может быть зафиксирована многократно.

\subsubsection{Понятия предметной области}

Обращение~--- зафиксированная жителем неисправность или дефект городской инфраструктуры (яма на дороге, неисправный фонарь, переполненный мусорный контейнер и~т.\,п.), требующий реагирования со стороны ответственной городской службы. Обращение содержит текстовое описание, фотографию, категорию проблемы и географические координаты.

Категория проблемы~--- классификационный признак, позволяющий отнести обращение к одному из тематических направлений: дороги, освещение, благоустройство, ЖКХ, экология и другие. Категория определяет, какая городская служба является ответственной за устранение данного типа проблемы.

Ответственная служба~--- организация или подразделение, в компетенцию которого входит устранение проблем определённой категории. Для каждой службы в системе хранятся контактные данные: телефон, адрес, режим работы.

Статус обращения~--- текущее состояние обращения в системе. Жизненный цикл обращения включает следующие состояния: <<новое>> (обращение создано жителем), <<принято>> (диспетчер подтвердил обращение), <<в работе>> (обращение передано ответственной службе), <<решено>> (проблема устранена), <<отклонено>> (обращение признано нерелевантным).

Дубликат обращения~--- повторное обращение, описывающее ту же проблему, что и ранее поданное, и расположенное в непосредственной географической близости от него (в пределах заданного радиуса). Механизм обнаружения дубликатов реализован на стороне сервера.

Диспетчер~--- оператор системы, осуществляющий обработку поступивших обращений через веб-панель: просмотр на карте, определение ответственной службы, изменение статуса и контроль хода устранения проблемы.

Диспетчеризация~--- процесс обработки обращений, включающий их приём, классификацию, определение ответственной службы, изменение статуса и контроль результата.

\subsubsection{Моделирование процессов}

Текущий процесс фиксации и диспетчеризации городских проблем представлен на контекстной диаграмме (рисунок~\ref{fig:current_a0}). Процесс предполагает участие трёх механизмов: диспетчера, программной системы фиксации событий и сотрудников служб. На вход поступает проблема инфраструктуры, управление осуществляется на основании законов РФ и справочника категорий проблем и служб.

\begin{figure}[H]
    \centering
    \includegraphics[width=0.8\linewidth]{to-be_a-0.png}
    \caption{Текущий процесс фиксации и диспетчеризации городских проблем (контекстная диаграмма, уровень А-0)}
    \label{fig:current_a0}
\end{figure}

Декомпозиция текущего процесса (рисунок~\ref{fig:current_decomp}) включает четыре подпроцесса:
\begin{enumerate}
    \item фиксация проблемы жителем~--- житель фиксирует проблему, на выходе формируются данные об обращении;
    \item проверка на дубликаты~--- система выполняет геопространственный поиск существующих обращений в радиусе 50~м;
    \item диспетчеризация~--- диспетчер обрабатывает обращение: просматривает, определяет ответственную службу, изменяет статус;
    \item контроль исполнения~--- отслеживание статуса обращения, фиксация результата, информирование жителя.
\end{enumerate}

\begin{figure}[H]
    \centering
    \includegraphics[width=0.8\linewidth]{to_be_a0.png}
    \caption{Декомпозиция текущего процесса (уровень А0)}
    \label{fig:current_decomp}
\end{figure}

Подпроцесс <<Диспетчеризация>> детализирован на диаграмме в нотации IDEF3 (рисунок~\ref{fig:dispatcher_idef3}). Диаграмма описывает последовательность действий оператора при обработке обращения: просмотр обращения, принятие решения о корректности, определение ответственной службы, передача в работу или отклонение, фиксация действий в журнале аудита.

\begin{figure}[H]
    \centering
    \includegraphics[width=0.7\linewidth]{to_be_dispatcher_idef3_decompose.png}
    \caption{Декомпозиция подпроцесса <<Диспетчеризация>> в нотации IDEF3}
    \label{fig:dispatcher_idef3}
\end{figure}

Планируемый процесс с учётом мобильного приложения представлен на контекстной диаграмме (рисунок~\ref{fig:planned_a0}). Ключевым отличием от текущего процесса является появление жителя в качестве полноценного участника, взаимодействующего с программной системой посредством мобильного приложения.

\begin{figure}[H]
    \centering
    \includegraphics[width=0.8\linewidth]{current_a-0.png}
    \caption{Планируемый процесс с учётом мобильного приложения (контекстная диаграмма, уровень А-0)}
    \label{fig:planned_a0}
\end{figure}

Декомпозиция планируемого процесса (рисунок~\ref{fig:planned_decomp}) включает те же четыре подпроцесса, однако подпроцесс <<Фиксация проблемы жителем>> теперь выполняется через мобильное приложение, обеспечивающее автоматическое определение координат, выбор категории проблемы из справочника, прикрепление фотографий с камеры устройства и передачу обращения на сервер.

\begin{figure}[H]
    \centering
    \includegraphics[width=0.8\linewidth]{current_a0.png}
    \caption{Декомпозиция планируемого процесса с учётом мобильного приложения (уровень А0)}
    \label{fig:planned_decomp}
\end{figure}

\subsection{Обзор аналогов}

Для обоснования разработки собственной системы были рассмотрены существующие решения в области фиксации городских проблем и обработки обращений граждан.

\textbf{Платформа обратной связи <<Госуслуги. Решаем вместе>>} (рисунок~\ref{fig:analog-gosuslugi})~--- федеральный сервис, интегрированный с порталом Госуслуг. Достоинства: федеральный охват, включая город Киров; наличие мобильного приложения; поддержка фотографий и статусов. Недостатки: нет механизма дедупликации; автоматическое определение GPS-координат не является приоритетным сценарием; диспетчерская панель с картографическим представлением отсутствует.

\begin{figure}[H]
    \centering
    \includegraphics[width=0.5\linewidth]{gosuslugi-reshaem-vmeste.jpg}
    \caption{Мобильное приложение <<Госуслуги. Решаем вместе>>}
    \label{fig:analog-gosuslugi}
\end{figure}

\textbf{Портал <<Наш город>> (Москва)} (рисунок~\ref{fig:analog-nashgorod})~--- портал для подачи жалоб на состояние городской инфраструктуры с развитой системой категоризации, механизмом верификации и мобильным приложением. Недостатки: доступен исключительно для жителей Москвы; автоматическая дедупликация по радиусу отсутствует.

\begin{figure}[H]
    \centering
    \includegraphics[width=0.5\linewidth]{nash-gorod-moskva-user.jpg}
    \caption{Мобильное приложение <<Наш город>> (Москва)}
    \label{fig:analog-nashgorod}
\end{figure}

\textbf{<<Добродел>> (Московская область)} (рисунок~\ref{fig:analog-dobrodel})~--- портал обращений жителей Московской области с механизмом <<коллективного сообщения>> и кнопкой <<Поддержать>>. Мобильное приложение поддерживает фотофиксацию, автоопределение геолокации и push-уведомления. Недостатки: ограничен территорией Московской области; автоматическая дедупликация по радиусу не реализована; интерфейс перегружен дополнительными функциями.

\begin{figure}[H]
    \centering
    \includegraphics[width=0.5\linewidth]{dobrodel-user.jpg}
    \caption{Мобильное приложение <<Добродел>>}
    \label{fig:analog-dobrodel}
\end{figure}

\textbf{<<Наш Санкт-Петербург>>} (рисунок~\ref{fig:analog-nashspb})~--- региональный сервис для фиксации городских проблем в Санкт-Петербурге. Мобильное приложение поддерживает фотофиксацию, указание местоположения, push-уведомления. Недостатки: ограничен территорией Санкт-Петербурга; автоматическая дедупликация не реализована; подтверждение устранения заявителем не предусмотрено.

\begin{figure}[H]
    \centering
    \includegraphics[width=0.5\linewidth]{nash-spb-user.jpg}
    \caption{Мобильное приложение <<Наш Санкт-Петербург>>}
    \label{fig:analog-nashspb}
\end{figure}

\textbf{FixMyStreet} (рисунок~\ref{fig:analog-fixmystreet})~--- британская платформа с открытым исходным кодом. Достоинства: автоматическое определение геолокации, интерактивная карта, автоматическая маршрутизация в муниципальный орган, открытый исходный код. Недостатки: ориентирован на административное деление Великобритании; push-уведомления не реализованы; адаптация под российские реалии не выполнена.

\begin{figure}[H]
    \centering
    \includegraphics[width=0.5\linewidth]{fixmystreet-user.png}
    \caption{Мобильное приложение FixMyStreet}
    \label{fig:analog-fixmystreet}
\end{figure}

Проведённый анализ выявил, что среди рассмотренных аналогов нет решения, применимого к городу Кирову и обеспечивающего полную связку <<фиксация~--- диспетчеризация на карте~--- маршрутизация к ответственной службе~--- дедупликация по радиусу>>. Ни одно из решений не реализует автоматическую дедупликацию обращений по заданному радиусу. Региональные решения (<<Наш город>>, <<Добродел>>, <<Наш Санкт-Петербург>>) обладают развитыми мобильными приложениями, однако территориально ограничены. Таким образом, разработка собственной системы обоснована.
